\problemname{Lawn Stripes}

Lawn stripes are drawn using a roller to flatten grass blades in a particular
direction and impart a sheen that, at a distance, is distinguishable from
blades flattened in another direction. This is typically for aesthetic appeal.

We are going to draw lawn stripes in our square garden using a roller in any
of the $4$ cardinal directions. The roller must stay on the lawn throughout.
The final direction of a piece of lawn is determined by the
last time the roller ran over it, including starting from or finishing there.

\begin{figure}[!h]
\centering
\includegraphics[height=0.4\textwidth]{sample1.pdf}
\caption{Illustration of Sample Input 1.}
\end{figure}

Given a description of the lawn and the starting point of the roller, determine
how to draw the pattern.

\section*{Input}

\begin{itemize}
  \item The first line of input contains the width and depth in metres of the lawn, $n$ ($1 \le n \le 100$)

  \item The second line of input contains the row and column of the roller, $r$ and $c$ ($1 \le r, c \le n$).

  \item Each of the next $n$ lines contains $n$ characters from \texttt{(l|r|u|d)} giving the desired
        direction of a square metre of lawn.
\end{itemize}

\section*{Output}

If it is not possible to draw the pattern, output \texttt{impossible}. Otherwise, demonstrate a solution in
under $10n^2$ moves by outputting the number of moves to use on one line, followed by the moves on the remaining lines.

Each move consists of a distance in metres $d$ and a direction from \texttt{(l|r|u|d)}. A roller
following the instructions must roll over every square metre without leaving the lawn.
